\documentclass{article}
\usepackage[utf8]{inputenc}
\usepackage[T1]{fontenc}
\usepackage{geometry}
\usepackage{graphicx}
\usepackage{float}
\usepackage{hyperref}
\usepackage[english,indonesian]{babel}

\title{Optimasi Pemrosesan Template LaTeX Menggunakan Pendekatan Berbasis AI}
\author{Academic Support Agent}
\date{\today}

\begin{document}
\maketitle

\begin{abstract}

\end{abstract}

\chapter{PENDAHULUAN}


\chapter{TINJAUAN PUSTAKA}


% Original content (commented out):
% %%%% TU Ilmenau - Fachgebiet Fahrzeugtechnik %%%%
%%% ----- Vorlage Bericht Fachpraktikum ----- %%%
%%%%%%%%%%%%%%%%%%%%%%%%%%%%%%%%%%%%%%%%%%%%%%%%%
% Version: 0.1                                  %
% erstellt am 05.03.2023                        %
% von kathleen.brueckner@tu-ilmenau.de          %
%%%%%%%%%%%%%%%%%%%%%%%%%%%%%%%%%%%%%%%%%%%%%%%%%
%%%%%%%%%% File-Setup %%%%%%%%%%

\documentclass[
a4paper,
oneside,
12pt,
titlepage]
{scrartcl}

% Layout 
\usepackage{geometry}
\geometry{
	top=2.5cm, % oberer Rand
	bottom=2.5cm, % unterer Rand
	left=3cm, % linker R.
	right=2cm, % rechter R.
	headheight=0.75cm, % Header height
	footskip=1.5cm, % Space from the bottom margin to the baseline of the footer
	headsep=0.75cm, % Space from the top margin to the baseline of the header
	%showframe, % Auskommentieren, wenn Rahmen des Textbereichs angezeigt werden soll
}


% Text
% Wenn als Compiler pdfLatex genutzt werden soll
\usepackage[T1]{fontenc}
\usepackage[utf8]{inputenc}

%\usepackage{fontspec} % nur wenn als Compiler: LuaLatex gewählt wird
\usepackage[ngerman]{babel}
\usepackage[onehalfspacing]{setspace} % Zeilenabstand 1,5 fach
\usepackage[skip = 9pt, indent= 0pt]{parskip} % Absatzformatierung: Abstand zwischen Absätzen, Einrücken Beginn

% Pakete Mathe/Formeln
\usepackage{amsmath}
\usepackage{amsfonts, amsthm, amssymb}
\usepackage{braket, nicefrac}

\usepackage{siunitx} % Nutzung SI Einheiten
\usepackage{enumitem} % Anpassug Standard Umgebung von description, itemize und enumerate
\usepackage{multicol} % Nutzung von Spalten

% Einfügen von Bildern
\usepackage{graphicx, float} 
\usepackage{pgfplots}\usepgfplotslibrary{units}
\pgfplotsset{width=10cm,compat=1.9}
\usepackage{caption}

% Pakete zum zeichnen
\usepackage{tikz}
\usepackage{tkz-euclide}
\usetikzlibrary{math}
\usepackage{fourier}
\usetikzlibrary{quotes,angles}
\usepackage{tkz-euclide}
\usetikzlibrary{calc}
\usepackage[german]{csquotes}

% Hyperlinks
\usepackage[colorlinks=true, linkcolor=blue, citecolor=blue, urlcolor=blue, bookmarks=true, breaklinks=true]{hyperref}
\labelformat{equation}{(#1)} 
\def\equationautorefname#1{}%for autoref, gobble the space

%%%%%%%%%% Ende File-Setup

%%%%%%%%%% Literaturverzeichnis 
\usepackage[style=ieee]{biblatex}
\addbibresource{references.bib}
%%%%%%%%%%

% relativer Ordnerpfad zu Bildern
\graphicspath{{graphics/}{bilder/}{./}}


\begin{document}

% Titelseite
\begin{titlepage}
	\centering
	\includegraphics[width=0.8\textwidth]{logo.png}\par\vspace{1cm}
	{\scshape\LARGE Praktikumsbericht \par} 
	\vspace{1cm}
	{\scshape\Large Fachpraktikum\par}
	\vspace{1.5cm}
	{\huge\bfseries Unternehmen/Forschungsinstitut\par}
        {\LARGE\bfseries Thema\par}
	\vspace{2cm}
	{\Large\itshape Vorname Nachname}\par
	\texttt{vorname.nachname@tu-ilmenau.de}\par
        \texttt{Matrikelnummer: 10815}\par
	\vspace{0.5cm}
	\vfill
	Praktikumsbetreuung:\par 
        Name betriebl. Betreuer,M.Sc /Name Betreuer FG, M.Sc
	\vfill
	{\large \today\par} % heutige Datum
\end{titlepage}

\tableofcontents % Inhaltsverzeichnis anlegen
\newpage

% Beginn Text
Der Praktikumsbericht ist in einem wissenschaftlichen Schreibstil anzufertigen, dabei können bildliche Darstellungen, Skizzen, Diagramme und so weiter nützlich sein. Die Regeln richtigen Zitierens nach DIN ISO 690 sind zu beachten. Außerdem sollte der Bericht in dieser Formatvorlage erstellt werden.
%\cite{anarticle, thesis, report, conv, unpub, notpub, notpub}.

Um den inhaltlichen Rahmen des Berichts abzustecken, ist hier ein Vorschlag über die Aufteilung der Kapitel mit den erforderlichen Inhalten aufgezeigt. 

\section{Einleitung}

In der Einleitung kann eine Übersicht über das Praktikum gegeben werden sowie eine kurze Beschreibung des Unternehmens. Außerdem sollte auf die Motivation für das Praktikum eingegangen werden, sowie erläutert werden, welche persönlichen Erwartungen an die Praktikumsstelle gestellt sind.

\section{Tätigkeiten im Fachpraktikum}
In diesem Kapitel sollte vorgestellt werden, welche Tätigkeiten und Aufgaben im Fachpraktikum erfüllt worden sind. Dabei ist herauszustellen welche Inhalte aus dem Studium relevant waren oder ob beispielsweise theoretisches Wissen eingebracht/gefestigt werden konnte. 

\subsection{Unterkapitel 1}
Hier kann eine Gesamtübersicht über die fachliche und zeitliche Gliederung gegeben werden. Sollte das Praktikum in verschiedenen Abteilungen durchgeführt worden sein, sind diese vorzustellen. Die Tätigkeitsbereiche und Arbeitsbedingungen können aufgezeigt werden. Es soll nicht jede einzelne Tätigkeit im Detail beschrieben werden. Vielmehr ist sich hier auf das Wesentliche zu beschränken. Die Vorgehensweise zur Lösung einer Aufgabe oder eines Problems ist von besonderem Interesse. Es ist möglich zu beschreiben, was tägliche Aufgaben waren und wie diese von Nutzen sind. Eine Aufgabe oder Projekt, welches über längere Zeit bearbeitet wurde, sollte vorgestellt werden. Es ist auch möglich einen Tages- oder Wochenbericht aufzuzeigen.


\subsection{Unterkapitel 2}
Im Text wird auf Abbildungen, Gleichungen und Tabellen referenziert.

\begin{equation}\label{eq:beispiel}
p_e = \frac{M_D \cdot 2\pi}{V_h \cdot i}
\end{equation}

Um die Vorlage besser nutzen zu können sind hierfür bereits Beispiele wie Gleichung \ref{eq:beispiel} eingefügt.\cite{abook}


% Grafik einfügen
\begin{figure}
    \centering
    \begin{tikzpicture}
        \node[] (chamber) at (0,0) 
          {\includegraphics[width=0.7\linewidth]{MPS Bild.jpg}};
          \draw[thick, white] (-0.5,-0.1) circle (1.5cm);
          \draw[<-,white] (-2,0) -- (-3.5,1) node[fill=white,text=black]{Motor};
          \draw[thick, white] (2.5,0.5) circle (0.5cm);
          \draw[<-,white] (2.5,1) -- (2.5,2.5) node[fill=white,text=black]{Steuergerät};
      \end{tikzpicture}
    \caption[MPS]{Motorprüfstand mit Versuchsträger}
    \label{fig:mps}
\end{figure}

 Hier ist ein hochdynamischer Motorprüfstand in Abbildung \ref{fig:mps} dargestellt. Beschriftungen in Grafiken können direkt mit tikz eingefügt werden. Da \LaTeX \  ein Textsatzsystem ist, werden Grafiken automatisch an eine dem Textfluss angepasste Position verschoben. Gerade bei technischen Dokumenten ist dies kein Nachteil, sondern dient dem besseren Lesefluss, weil keine größeren Lücken entstehen. Bei Bedarf lässt sich die Positionierung jedoch auch erzwingen.
 
 \subsection{Unterkapitel 3}



 % Tabelle einfügen
\begin{table}[ht]
    \centering
    \begin{tabular}{llr}
\hline
\multicolumn{2}{c}{Objekte} \\ 
\cline{1-2}
Objekt    & Beschreibung & Länge (\unit{\metre}) \\
\hline
Objekt 1      & Beschreibung 1    & \qty{1.00}{\metre}      \\
Objekt 2      & Beschreibung 2    & \qty{2.00}{\metre}       \\ 
Objekt 3      & Beschreibung 3    & \qty{3.50}{\metre}      \\

\hline
\end{tabular}
    \caption{Beispieltabelle}
    \label{tab:bsp}
\end{table}

Wie Tabellen aussehen können, zeigt Tabelle \ref{tab:bsp}.

\section{Fazit}
Im Fazit ist eine Zusammenfassung über die gelernten Inhalte während des Praktikums möglich. Es kann aufgezeigt werden, ob theoretische Konzepte aus dem Studium im Praktikum angewendet werden konnte, was sich die Studierenden noch gewünscht hätten einbringen zu können oder welche Erkenntnisse gewonnen und welche Wissenslücken aufgezeigt wurden. Außerdem sollte reflektiert werden, inwieweit die eigenen Erwartungen an die Praktikumsstelle erfüllt oder nicht erfüllt werden konnten und welche Bedeutung das Praktikum für ihr Studium hat.

\newpage
\printbibliography

% Um CAD Zeichnungen, Skizzen oder data files, usw. in den Anhang zu bringen:
%\newpage
%\appendix
%\section{Anhang}


\end{document}

\end{document}